\documentclass[a4paper,12pt]{scrreprt}%Dokumentklasse
\usepackage[left=4cm, right=2cm, top=2.8cm, bottom=2cm,
	includehead=true, % Kopfzeile innerhalb des Textkörper, also nicht im Rand
    includefoot=false, % muss true gesetzt werden, wenn eine Fußzeile genutzt wird!!
    %showframe,%Anzeigen der Frames für das Messen
    bindingoffset=0mm]{geometry}
    
\renewcommand{\footnotesize}{\fontsize{10}{10} \selectfont}%schriftgröße der Fußziele  
    
\usepackage[onehalfspacing]{setspace}
% ============= Packages =============

% Dokumentinformationen
\usepackage[
	pdftitle={Analyse von Erhebungstechniken},
	pdfsubject={},
	pdfauthor={Ingo Sewcz},
	pdfkeywords={Erhebungstechniken, Dokumentenanalyse, Laufzettel, Interview},	
	%Links nicht einrahmen
	hidelinks
]{hyperref}


% Standard Packages
\usepackage[utf8]{inputenc}
\usepackage[ngerman]{babel}
\usepackage[T1]{fontenc}
\usepackage{graphicx, subfig}
\graphicspath{{img/}}
\usepackage[
  automark,
  autooneside=false,% <- needed if you want to use \leftmark and \rightmark in a onesided document
  headsepline
]{scrlayer-scrpage}
\usepackage{lmodern}
\usepackage{color}


% zusätzliche Schriftzeichen der American Mathematical Society
\usepackage{amsfonts}
\usepackage{amsmath}

%nicht einrücken nach Absatz
\setlength{\parindent}{0pt}
%Aber Abstand zwischen Absätzen
\setlength{\parskip}{12pt}


%Zitation
\usepackage[backend=biber, style=authoryear]{biblatex}
\usepackage[babel, german=quotes]{csquotes}
\setlength{\bibitemsep}{12pt}
\addbibresource{Literatur.bib}
%Zietert wird dann an der gewünschten Stelle wie folgt:
%\parencite[vgl.][S.60]{mueller2000}
%\cite[vgl.][S.60]{mueller2000}

% ============= Kopf- und Fußzeile =============
\clearpairofpagestyles
\ihead{\leftmark}%LINKS
\chead{}%MITTE
\ohead{\thepage}%RECHTS
\renewcommand*{\headfont}{\normalfont}
\pagestyle{scrheadings}
\renewcommand*{\chapterpagestyle}{scrheadings} 

\RedeclareSectionCommands[
  beforeskip=-.5\baselineskip,
  afterskip=.25\baselineskip
]{chapter,section,subsection,subsubsection}
\RedeclareSectionCommands[
  beforeskip=.5\baselineskip,
  afterskip=-1em]{paragraph,subparagraph}

% ============= Package Einstellungen & Sonstiges ============= 
%Besondere Trennungen
\hyphenation{De-zi-mal-tren-nung}


% ============= Dokumentbeginn =============

\begin{document}
%Seiten ohne Kopf- und Fußzeile sowie Seitenzahl
\pagestyle{empty}
%Seitennummerierung auf Römisch umstellen
\pagenumbering{Roman}

\begin{center}
\begin{tabular}{p{410pt}}%die Breite muss hier gesetzt sein, da sonst eine Overflow Warning in Latex auftritt.


\begin{center}
\includegraphics[scale=0.4]{img/fomlogo.png}
\end{center}


\\

\begin{center}
\LARGE{\textsc{
Serverlastreduzierung\\durch den Einsatz von Angular\\
}}
\end{center}

\\

\begin{center}
\textbf{\Large{Seminararbeit}}
\end{center}


\begin{center}
im Fach\\
IT-Infrastruktur \\
des Studiengangs Wirtschaftsinformatik - Business Information Systems \\
(Fachsemester III)
\end{center}


\begin{center}
vorgelegt von
\end{center}

\begin{center}
\large{\textbf{Ingo Sewcz}} \\
\small{geboren am 24.07.1994 in Marl}
\end{center}

\begin{center}
\large{am 12.01.2020, in Münster}
\end{center}

\\

\begin{center}
\begin{tabular}{lll}
\textbf{Erstprüfer:} & & Prof. Dr. Gregor Hülsken\\
\end{tabular}
\end{center}

\end{tabular}
\end{center}


%\addsec*{Danksagungen}
\label{danksagungen}
Weit hinten, hinter den Wortbergen, fern der Länder Vokalien und Konsonantien leben die Blindtexte. Abgeschieden wohnen Sie in Buchstabhausen an der Küste des Semantik, eines großen Sprachozeans. Ein kleines Bächlein namens Duden fließt durch ihren Ort und versorgt sie mit den nötigen Regelialien. Es ist ein paradiesmatisches Land, in dem einem gebratene Satzteile in den Mund fliegen. Nicht einmal von der allmächtigen Interpunktion werden die Blindtexte beherrscht – ein geradezu unorthographisches Leben. Eines Tages aber beschloß eine kleine Zeile Blindtext, ihr Name war Lorem Ipsum, hinaus zu gehen in die weite Grammatik. Der große Oxmox riet ihr davon ab, da.

\newpage
\addsec*{Sperrvermerk}
\label{sperrvermerk}
Weit hinten, hinter den Wortbergen, fern der Länder Vokalien und Konsonantien leben die Blindtexte. Abgeschieden wohnen Sie in Buchstabhausen an der Küste des Semantik, eines großen Sprachozeans. Ein kleines Bächlein namens Duden fließt durch ihren Ort und versorgt sie mit den nötigen Regelialien. Es ist ein paradiesmatisches Land, in dem einem gebratene Satzteile in den Mund fliegen. Nicht einmal von der allmächtigen Interpunktion werden die Blindtexte beherrscht – ein geradezu unorthographisches Leben. Eines Tages aber beschloß eine kleine Zeile Blindtext, ihr Name war Lorem Ipsum, hinaus zu gehen in die weite Grammatik. Der große Oxmox riet ihr davon ab, da.

%\addsec*{Zusammenfassung / Abstract}
\label{sec:zusammenfassung}


\minisec{Abstract}
\label{abstract}


% Beendet eine Seite und erzwingt auf den nachfolgenden Seiten die Ausgabe aller Gleitobjekte (z.B. Abbildungen), die bislang definiert, aber noch nicht ausgegeben wurden. Dieser Befehl fügt, falls nötig, eine leere Seite ein, sodaß die nächste Seite nach den Gleitobjekten eine ungerade Seitennummer hat. 
\cleardoubleoddpage

% pagestyle für gesamtes Dokument aktivieren
\pagestyle{scrheadings}

%Inhaltsverzeichnis
\tableofcontents

%Verzeichnis aller Bilder
\newpage
\addcontentsline{toc}{section}{Abbildungsverzeichnis}
\listoffigures

%Verzeichnis aller Tabellen
\newpage
\addcontentsline{toc}{section}{Tabellenverzeichnis}
\listoftables

\newpage
%Seitennummerierung auf Arabisch umstellen
\pagenumbering{arabic}

\chapter{Einleitung}
\label{sec:einleitung}
Weit hinten, hinter den Wortbergen, fern der Länder Vokalien und Konsonantien leben die Blindtexte. Abgeschieden wohnen Sie in Buchstabhausen an der Küste des Semantik, eines großen Sprachozeans. Ein kleines Bächlein namens Duden fließt durch ihren Ort und versorgt sie mit den nötigen Regelialien. Es ist ein paradiesmatisches Land, in dem einem gebratene Satzteile in den Mund fliegen. Nicht einmal von der allmächtigen Interpunktion werden die Blindtexte beherrscht – ein geradezu unorthographisches Leben. Eines Tages aber beschloß eine kleine Zeile Blindtext, ihr Name war Lorem Ipsum, hinaus zu gehen in die weite Grammatik. Der große Oxmox riet ihr davon ab, da.

\chapter{Grundlagen}
\label{sec:grundlagen}


\section{Beispielkapitel}
\label{sec:beispiel}
\begin{figure}[htb]
  \centering  
  \includegraphics[scale=0.2]{img/fomlogo.png}
  \caption{Star Wars Logo}
  \label{fig:starwars}
\end{figure}
Weit hinten, hinter den Wortbergen, fern der Länder Vokalien und Konsonantien leben die Blindtexte (siehe Abb \ref{fig:starwars}). Abgeschieden wohnen Sie in Buchstabhausen an der Küste des Semantik, eines großen Sprachozeans. Ein kleines Bächlein namens Duden fließt durch ihren Ort und versorgt sie mit den nötigen Regelialien. Es ist ein paradiesmatisches Land, in dem einem gebratene Satzteile in den Mund fliegen. Nicht einmal von der allmächtigen Interpunktion werden die Blindtexte beherrscht – ein geradezu unorthographisches Leben. Eines Tages aber beschloß eine kleine Zeile Blindtext, ihr Name war Lorem Ipsum, hinaus zu gehen in die weite Grammatik. Der große Oxmox riet ihr davon ab, da\footnote{\cite[vgl.][S.60]{Khan2016}}.
Weit hinten, hinter den Wortbergen, fern der Länder Vokalien und Konsonantien leben die Blindtexte (siehe Abb \ref{fig:starwars}). Abgeschieden wohnen Sie in Buchstabhausen an der Küste des Semantik, eines großen Sprachozeans\footnote{Test einer Fußzeile}. Ein kleines Bächlein namens Duden fließt durch ihren Ort und versorgt sie mit den nötigen Regelialien.\footnote{und noch eine Zweite um den Abstand zu messen} Es ist ein paradiesmatisches Land, in dem einem gebratene Satzteile in den Mund fliegen. Nicht einmal von der allmächtigen Interpunktion werden die Blindtexte beherrscht – ein geradezu unorthographisches Leben. Eines Tages aber beschloß eine kleine Zeile Blindtext, ihr Name war Lorem Ipsum, hinaus zu gehen in die weite Grammatik.

Weit hinten, hinter den Wortbergen, fern der Länder Vokalien und Konsonantien leben die Blindtexte (siehe Abb \ref{fig:starwars}). Abgeschieden wohnen Sie in Buchstabhausen an der Küste des Semantik, eines großen Sprachozeans. Ein kleines Bächlein namens Duden fließt durch ihren Ort und versorgt sie mit den nötigen Regelialien. Es ist ein paradiesmatisches Land, in dem einem gebratene Satzteile in den Mund fliegen. Nicht einmal von der allmächtigen Interpunktion werden die Blindtexte beherrscht – ein geradezu unorthographisches Leben. Eines Tages aber beschloß eine kleine Zeile Blindtext, ihr Name war Lorem Ipsum, hinaus zu gehen in die weite Grammatik.
Weit hinten, hinter den Wortbergen, fern der Länder Vokalien und Konsonantien leben die Blindtexte (siehe Abb \ref{fig:starwars}). Abgeschieden wohnen Sie in Buchstabhausen an der Küste des Semantik, eines großen Sprachozeans. Ein kleines Bächlein namens Duden fließt durch ihren Ort und versorgt sie mit den nötigen Regelialien. Es ist ein paradiesmatisches Land, in dem einem gebratene Satzteile in den Mund fliegen. Nicht einmal von der allmächtigen Interpunktion werden die Blindtexte beherrscht – ein geradezu unorthographisches Leben. Eines Tages aber beschloß eine kleine Zeile Blindtext, ihr Name war Lorem Ipsum, hinaus zu gehen in die weite Grammatik.

Weit hinten, hinter den Wortbergen, fern der Länder Vokalien und Konsonantien leben die Blindtexte (siehe Abb \ref{fig:starwars}). Abgeschieden wohnen Sie in Buchstabhausen an der Küste des Semantik, eines großen Sprachozeans. Ein kleines Bächlein namens Duden fließt durch ihren Ort und versorgt sie mit den nötigen Regelialien. Es ist ein paradiesmatisches Land, in dem einem gebratene Satzteile in den Mund fliegen. Nicht einmal von der allmächtigen Interpunktion werden die Blindtexte beherrscht – ein geradezu unorthographisches Leben. Eines Tages aber beschloß eine kleine Zeile Blindtext, ihr Name war Lorem Ipsum, hinaus zu gehen in die weite Grammatik.
Weit hinten, hinter den Wortbergen, fern der Länder Vokalien und Konsonantien leben die Blindtexte (siehe Abb \ref{fig:starwars}). Abgeschieden wohnen Sie in Buchstabhausen an der Küste des Semantik, eines großen Sprachozeans. Ein kleines Bächlein namens Duden fließt durch ihren Ort und versorgt sie mit den nötigen Regelialien. Es ist ein paradiesmatisches Land, in dem einem gebratene Satzteile in den Mund fliegen. Nicht einmal von der allmächtigen Interpunktion werden die Blindtexte beherrscht – ein geradezu unorthographisches Leben. Eines Tages aber beschloß eine kleine Zeile Blindtext, ihr Name war Lorem Ipsum, hinaus zu gehen in die weite Grammatik.

Weit hinten, hinter den Wortbergen, fern der Länder Vokalien und Konsonantien leben die Blindtexte (siehe Abb \ref{fig:starwars}). Abgeschieden wohnen Sie in Buchstabhausen an der Küste des Semantik, eines großen Sprachozeans. Ein kleines Bächlein namens Duden fließt durch ihren Ort und versorgt sie mit den nötigen Regelialien. Es ist ein paradiesmatisches Land, in dem einem gebratene Satzteile in den Mund fliegen. Nicht einmal von der allmächtigen Interpunktion werden die Blindtexte beherrscht – ein geradezu unorthographisches Leben. Eines Tages aber beschloß eine kleine Zeile Blindtext, ihr Name war Lorem Ipsum, hinaus zu gehen in die weite Grammatik.
Weit hinten, hinter den Wortbergen, fern der Länder Vokalien und Konsonantien leben die Blindtexte (siehe Abb \ref{fig:starwars}). Abgeschieden wohnen Sie in Buchstabhausen an der Küste des Semantik, eines großen Sprachozeans. Ein kleines Bächlein namens Duden fließt durch ihren Ort und versorgt sie mit den nötigen Regelialien. Es ist ein paradiesmatisches Land, in dem einem gebratene Satzteile in den Mund fliegen. Nicht einmal von der allmächtigen Interpunktion werden die Blindtexte beherrscht – ein geradezu unorthographisches Leben. Eines Tages aber beschloß eine kleine Zeile Blindtext, ihr Name war Lorem Ipsum, hinaus zu gehen in die weite Grammatik.


\section{Beispielunterkapitel}
\label{subsec:beispiel}

Weit hinten, hinter den Wortbergen, fern der Länder Vokalien und Konsonantien leben die Blindtexte (siehe Abb \ref{fig:starwars}). Abgeschieden wohnen Sie in Buchstabhausen an der Küste des Semantik, eines großen Sprachozeans. Ein kleines Bächlein namens Duden fließt durch ihren Ort und versorgt sie mit den nötigen Regelialien. Es ist ein paradiesmatisches Land, in dem einem gebratene Satzteile in den Mund fliegen. Nicht einmal von der allmächtigen Interpunktion werden die Blindtexte beherrscht – ein geradezu unorthographisches Leben. Eines Tages aber beschloß eine kleine Zeile Blindtext, ihr Name war Lorem Ipsum, hinaus zu gehen in die weite Grammatik.
Weit hinten, hinter den Wortbergen, fern der Länder Vokalien und Konsonantien leben die Blindtexte (siehe Abb \ref{fig:starwars}). Abgeschieden wohnen Sie in Buchstabhausen an der Küste des Semantik, eines großen Sprachozeans. Ein kleines Bächlein namens Duden fließt durch ihren Ort und versorgt sie mit den nötigen Regelialien. Es ist ein paradiesmatisches Land, in dem einem gebratene Satzteile in den Mund fliegen. Nicht einmal von der allmächtigen Interpunktion werden die Blindtexte beherrscht – ein geradezu unorthographisches Leben. Eines Tages aber beschloß eine kleine Zeile Blindtext, ihr Name war Lorem Ipsum, hinaus zu gehen in die weite Grammatik.

\chapter{Stand der Technik}
\label{cha:stand_der_technik}

\section{Klassifizierung von Beispielen}
\label{sec:klassifizierung}

\section{Themenbezogene Veröffentlichungen}
\label{sec:themenbezogene_veroeffentlichungen}

\chapter{Entwicklung der Sensorik}
\label{chap:entwicklung}


\section{Grobkonzept der Sensorik}
\label{sec:grobkonzept}

\chapter{Ergebnisse}
\label{sec:ergebnisse}


\chapter{Diskussion}
\label{sec:diskussion}

\section{Zusammenfassende Bewertung}
\label{sec:bewertung}

\section{Ausblick}
\label{sec:ausblick}

%Literaturverzeichnis
\newpage
%\bibliographystyle{unsrt}
%\bibliography{Literatur}
%Ersetzt duch:
\addcontentsline{toc}{section}{Literatur}
\printbibliography 

\newpage
%Seiten ohne Kopf- und Fußzeile sowie Seitenzahl
\pagestyle{empty}
\addsec*{Eidesstattliche Erklärung}
\label{erklaerung}

Hiermit versichere ich, die vorliegende Abschlussarbeit selbstständig und nur unter Verwendung der von mir angegebenen Quellen und Hilfsmittel verfasst zu haben. Sowohl inhaltlich als auch wörtlich entnommene Inhalte wurden als solche kenntlich gemacht. Die Arbeit hat in dieser oder vergleichbarer Form noch keinem anderem Prüfungsgremium vorgelegen.

\begin{center}
\begin{tabular}{p{410pt}}
\\[1.5cm]
Datum: \hrulefill\enspace Unterschrift: \hrulefill
\\[3.5cm]
\end{tabular}
\end{center}

\end{document}
